\section{最大熵原理：从约束到分布}
\label{sec:07-Information-Theory/01-Information-and-Entropy/04}

一个工程师要仿真零件寿命。她只知道两条硬信息：寿命非负，样本均值五年。仅凭这两条，该画什么分布？

画条钟形曲线？不行——Gaussian 伸到负半轴，零件不可能负寿命。画个三角分布？那对称性从哪儿来的？画个 [0, 10] 的均匀分布？凭什么上限是 10 而不是 15？

每条曲线都偷偷在数据里塞了未经验证的结构。问题不是``找不到分布''，而是``能找到太多分布''——需要一条原则来做选择，并且公开承认选了哪条原则。

\subsection{不是预测现实，是拒绝编造}

最大熵原理的想法并不复杂：

\begin{quote}
在满足已知约束的所有分布中，挑熵最大的那个——熵最大不是因为它``更真实''，而是因为它拒绝添加任何约束之外的结构。
\end{quote}

换个角度：已知约束划出了一片可行的分布集合。在这个集合里，任何一个比最大熵解更集中的分布，必然在某个维度上添加了你没有证据的偏好。最大熵解是这片集合里的``最不额外假设''的解——它把所有未被约束固定的概率质量，尽可能均匀地摊开。

这个原理是一种建模纪律，不是物理定律。没有理由认为大自然会主动最大化熵。但当你必须基于有限信息做出一个分布选择时，最大熵至少让你把额外假设的数量压到最低。

\subsection{第一道例题：只知道可能取值个数}

设 $X$ 有 $m$ 个可能值，唯一约束是归一化 $\sum_i p_i=1$。构造 Lagrange 函数：

\[
\mathcal L=-\sum_i p_i\log p_i+\lambda\Bigl(\sum_i p_i-1\Bigr).
\]

对 $p_i$ 求偏导（对数取自然底；换底只改变 $\lambda$ 的数值，不影响解的形式）：

\[
-(\log p_i+1)+\lambda=0,
\]

所有 $p_i$ 相等，结合归一化得

\[
p_i=\frac{1}{m}.
\]

均匀分布是唯一解。这回答了答案，但要注意前提：有限支撑、等权计数测度、无其他约束。并不是``什么都不知道就一律等可能''这样一句口号——前提里已经写死了有限的 $m$ 个类别和每一类权重相同。这些前提一旦松动，答案就不再是均匀分布。

\subsection{第二道例题：非负且只知均值}

在 $x\ge 0$ 上最大微分熵，约束为

\[
\int f(x)\,dx=1,\qquad \int x f(x)\,dx=\mu.
\]

把约束用 Lagrange 乘子写入目标，对 $f$ 逐点变分：

\[
-(\log f(x)+1)+\lambda_0+\lambda_1 x=0,
\]

解出 $f(x)\propto e^{\lambda_1 x}$。积分收敛需要 $\lambda_1<0$，改记 $\lambda=-\lambda_1>0$：

\[
f(x)\propto e^{-\lambda x},
\]

再用归一化和均值条件确定参数：

\[
f(x)=\frac{1}{\mu}e^{-x/\mu}.
\]

所以在非负支撑加固定均值的约束下，Exponential 分布最大熵。

实用中可以把这当作起点。如果目前只可靠地知道``等待时间非负''和``平均等待五分钟''，Exponential 模型没有擅自加入高峰时段、批量到达或固定检修周期等结构。它是一条可审计的基线——你清楚地知道模型所做的和没做的是什么。

但反过来，这不等于从均值``推导''出了无记忆性。午餐时段的集中到达、预约制度或零件磨损都会产生 Exponential 无法捕获的额外结构。一旦数据稳定地显示出这些现象，就应该把相应信息写进新约束或换用更合适的模型，而不是抱着``这已经是最大熵解''替原模型护航。

\subsection{第三道例题：知道均值和方差}

在全实线上固定

\[
\E[X]=\mu,\qquad \Var(X)=\sigma^2,
\]

最大微分熵解是 Gaussian：

\[
f(x)=\frac{1}{\sqrt{2\pi\sigma^2}}\exp\!\left[-\frac{(x-\mu)^2}{2\sigma^2}\right].
\]

用 KL 非负性证明很干脆。由 Jensen 不等式，$D(f\|g)\ge 0$。令 $g$ 为与 $f$ 同均值同方差的 Gaussian：

\[
D(f\|g)=-h(f)-\int f(x)\log g(x)\,dx.
\]

$\log g(x)$ 是 $x$ 的二次函数，积分只触及 $f$ 的前两阶矩。而 $f$ 与 $g$ 的前两阶矩相同，交叉项可换成 $g$ 自身的：

\[
-\int f(x)\log g(x)\,dx=-\int g(x)\log g(x)\,dx=h(g).
\]

于是 $0\le D(f\|g)=h(g)-h(f)$，即 $h(f)\le h(g)$。Gaussian 在给定的均值和方差下微分熵最大。

\subsection{一般形式：线性约束导出指数族}

推得更一般些。若约束是一组线性期望式：

\[
\E[T_k(X)]=c_k,
\]

最大熵解通常具有形式

\[
p_\theta(x)\propto\exp\!\left(\sum_k\theta_k T_k(x)\right).
\]

这就是指数族的结构。Bernoulli、Poisson、Gaussian、Exponential 等数十个常见分布，都可以通过指定支撑和充分统计量约束，从最大熵推导出来。Lagrange 乘子与约束期望之间的对偶关系，连接着最大熵、配分函数、凸优化和统计力学——不是巧合，而是同一套数学在不同语境下的复现。

\subsection{坐标系的选择不是免费的}

连续情形的微分熵依赖坐标，因此``最大熵''结论也会随坐标改变。对参数做非线性重参数化——比如把标准差改成方差——``均匀''或``最大微分熵''的答案可能不一样。更稳健的做法是明确引入参考测度 $m$，最小化相对熵：

\[
\min_p D(p\|m)
\]

并满足约束。这时的 $m$ 把基准结构写进了公式里，坐标依赖被相对熵的变换性质吸收。

\subsection{约束本身也会出错}

最大熵只保证不在约束外额外假设——它不保证约束就够了。若只约束均值，而真实数据是多峰的、带周期性的、或有硬边界，最大熵模型可能严重误导。

一个负责任的使用姿势应该至少交代：

\begin{itemize}
\tightlist
\item 支撑集是什么；
\item 参考测度（连续情形）或先验计数测度（离散情形）；
\item 每一条约束的来源和物理解释；
\item 约束是否相互兼容、是否存在可行解；
\item 解是否存在且可归一化；
\item 有哪些重要结构当前没有被约束，但可能对应用场景很关键。
\end{itemize}

最大熵把建模假设集中且透明地写成了约束——这恰好逼着建模者公开承认自己的选择，而不是把主观猜测藏在一条看起来客观的曲线后面。
